%% APPENDIX E -- radio-frequency interference screening.
%% Owner: appendix-triage.  Carries \label{app:rfi}.
\section{Radio-frequency interference screening}
\label{app:rfi}

Interference is the first alternative explanation for an unattributed
crossing, so the screening this survey has, and the coverage it does not
have, are both set out here.

Millimetre interferometry is an unusually poor environment for terrestrial
interference, though not an empty one: ITU Radio Regulations Article~5
\citep{ITU2024} carries active allocations throughout 90--275\,GHz, and WRC-19 identified
parts of 275--450\,GHz for land-mobile and fixed use. The specific hazard is
the \RadarLo--\RadarHi\,GHz Earth exploration-satellite allocation used by
spaceborne cloud radar --- CloudSat and, since 2024, EarthCARE
\citep{CloudSat2002,EarthCARE2023} --- a
documented contaminant of ALMA Band~3; \RadarNWin{} of the \RadarNTot{}
searched windows overlaps it, carrying \RadarNCross{} threshold crossings and
\RadarNStage{} rank-passing events, so no result here rests on it. What
favours this band is that terrestrial transmitters are sparse compared with
centimetre wavelengths, that the array sits at 5000\,m on a site chosen
partly for radio quietness, and, decisively, that an interferometer is not a
total-power detector: a signal not in the far field at the phase centre does
not fringe-track but acquires a rapidly varying geometric phase across the
array, and is suppressed in the correlation rather than imaged.

\emph{What screening exists, stated with its coverage, because the coverage
is uneven.} The allocation test above and the sky- and
intermediate-frequency tests below are computed for every searched window
from the released catalogue and are reproducible from it, as is the
\NCtrl-position rank screen, which asks whether an excess is at the star or
everywhere in the field and so serves as an interference test as well as a
prioritisation. Recurrence covers only the crossings that have a covering
repeat block, and the visibility fit --- which asks whether an excess has
the phase signature of a source at the phase centre, as a near-field emitter
would not --- is a consistency check rather than a filter. The correlator
flagging is the observatory's own, replayed verbatim from each block's
delivered calibration script, but it is not recorded per window in the
release. \emph{No per-channel interference flag of our own exists, and none
can be recovered.}

Four measurements support the reasoning, and each is a test the population
could have failed.

\emph{First, fixed sky frequency.} Interference is a property of the
observatory and not of the star, so a terrestrial carrier should appear at
the same sky frequency in every window that covers it. Here
\RfiEventLead; \RfiOtherWin{} further windows, toward \RfiOtherStarsRange{}
other stars in \RfiOtherBlocksRange{} other blocks, cover \RfiFreqThem,
and \RfiSameChanClause.\RfiNotTestableClause{}

\emph{Second, clustering in sky frequency.} If interference dominated the
crossing population, the crossings would pile up at particular sky
frequencies. They do cluster --- but on the survey's own molecular lines.
The release gives a peak frequency for \AppMNCrossF{} of the \NCross{}
crossings; the other \CsNCrossNoFreq{} are \CsCrossNoFreqStar{}
carbon-monoxide crossings, for which it gives the velocity offset from the
line instead, which is why Table~\ref{tab:ledger} lists all \NCross{} while
this test runs over \AppMNCrossF. Of those, \AppMTubeObs{}
fall inside the mask's own $\pm\DispoHalfKms$\,km\,s$^{-1}$ velocity window
against \AppMTubeNull{} expected if each window's peak fell anywhere on that
window's own channel grid, $\pv{\AppMTubeP}$. The densest band,
\AppMBandLo--\AppMBandHi\,GHz, holds \AppMBandN{} crossings, of which
\AppMBandNAttr{} are within that velocity window of CO($2{\to}1$); and once
the line-attributed crossings are set aside, the \AppMUnattrN{} that remain
show no clustering at all ($\pv{\AppMUnattrP}$). \emph{The concentration is
the attribution, not a residue left over after it.}

\emph{Third, fixed intermediate frequency.} An artefact of the signal chain
would sit at a fixed position within the spectral window rather than at a
fixed sky frequency. This test is run over every window for which the
catalogue releases a peak frequency --- \IfPopNWin{} of the \NWindows{}, of
which only \IfPopNCross{} carry a threshold crossing --- because such an artefact would
occupy the same intermediate frequency in all of them and not only where the
trigger fires. It does not: the peaks' fractional positions within their own
windows have median \RfiFracMed{} and are uniform.

\emph{Fourth, a cross-target occupancy screen over every searched window and
not only the crossings.} A terrestrial carrier should recur at the same sky
frequency toward \emph{unrelated} stars. Over all \RfiOccNWin{} searched
windows toward \RfiOccNStars{} stars, with each window's peak randomised on
its own channel grid, the number of cross-star coincidences within
\RfiOccTol\,MHz is \RfiOccObs{} against
\RfiOccNull$\,\pm\,$\RfiOccNullSd{} expected ($\pv{\RfiOccP}$); among the
\RfiOccSigN{} windows whose peak exceeds the trigger it is \RfiOccSigObs{}
against \RfiOccSigNull$\,\pm\,$\RfiOccSigNullSd, i.e.\ \emph{below} the
null, and with the molecular-line mask applied \RfiOccMskObs{} against
\RfiOccMskNull$\,\pm\,$\RfiOccMskNullSd. There is no excess to explain. The
screen is keyed on the star rather than on the product directory, which
matters: a star's own repeat observations are not different targets, and
every persistent feature --- a real line most of all --- would otherwise
enter as a cross-target coincidence. \RfiOccMaskN{} of those \RfiOccSigN{}
windows sit on a masked transition (\RfiOccSpeciesList), which is the point:
a CO line falls at almost the same sky frequency toward every
nearby star, so cross-star frequency coincidence is exactly what real line
emission produces, and the screen cannot distinguish it from interference
unless the mask is applied first.

\emph{One execution block fails these tests, and it fails them as a block
rather than as a carrier.} Its \DqBlockNCross{} crossings toward
\LgVerRejectStar{} are the highest statistics in the survey outside
$\beta$~Pictoris and HD~48370, and the reason is in the field and not at the
star: all \NCtrl{} control positions exceed the trigger, the ring's own
median running \DqBlockRingLo--\DqBlockRingHi{} against a survey median of
\DqSurvMed, so the window contains no null against which anything could be
localised. The visibility fit returns a large imaginary part, placing the
emission away from the phase centre; \DqBlockIfMult{} of the
\DqBlockNCross{} sit at the identical channel of two different spectral
windows, which is the signature of a fixed intermediate frequency; and the
block's \BdNSibling{} sibling executions of the same field, with the same
array and the same spectral setup, produce no crossing at all. It is
\emph{not} a terrestrial carrier at a fixed sky frequency: each of the four
frequencies is covered by \DqOtherWinLo--\DqOtherWinHi{} further windows
toward \DqOtherStarLo--\DqOtherStarHi{} other stars, and \DqOtherSame{} of
them carries a crossing in the same channel. \textbf{The evidence supports a
data-quality failure of one execution block, and that is how these four
crossings are dispositioned} (\S\ref{sec:bdblock}): under a rule fixed
before it was applied --- a window is flagged when the median of its own
control ensemble exceeds \DqFactor{} the survey median of that quantity ---
\DqNWin{} of the \NWindows{} windows are flagged, toward \DqNStar{} stars
(\DqStars), the remaining one being a resolved disc rather than a defect.
The flag is a statement about the window, not an attribution of the
feature.

\emph{None of this is a substitute for the ON/OFF cadence a dedicated search
uses, and we do not present it as one.} Accordingly we claim only that
\emph{no identified interference mechanism accounts for the unattributed
events}:
allocation tables cannot exclude harmonics, intermodulation products,
unconventional emitters or near-field sources, and an archival dataset
without a simultaneous off-source strategy cannot do better. What the four
tests establish is that the crossing population as a whole carries none of
the signatures interference leaves --- fixed sky frequency, fixed
intermediate frequency, or presence in unrelated targets --- which is a
statement about the population and not a clearance of any individual
crossing.

%% ------------------------------------------------------------------
%% NOTE: per 06_discussion.tex's own R6, this appendix no longer ends on
%% "the null result resting on internal validation alone" -- that statement
%% now lives once, in the discussion's lessons subsection.
%% REQUEST (appendix-triage -> integrator):
%%   fig:appm (the power comparison of the per-channel multiplicity statistic
%%   against the scan statistic) is deleted: it exists to compare our own two
%%   statistics, and the paper now states only the one it uses.  The
%%   \AppMPowerOld* and \RfiOcc{Prior,Dir,Stale}* macro families are no
%%   longer referenced and will be retired.  figures/appm_power.pdf is no
%%   longer included.
%% ------------------------------------------------------------------
